% ==============================================================================
% IEEE Sensors Journal (Feb 2026)
% Section IV-C4: Computational Efficiency, Stage Breakdown, and Parallel Scaling
% Section V: Conclusions and Future Prospects
% ==============================================================================

\subsection{Computational Efficiency, Stage Breakdown, and Multi-Core Scaling}
\label{sec:runtime_analysis}

Real-time or near-real-time throughput is an essential operational requirement for practical infrared focal plane array (IRFPA) sensing applications, including aerospace surveillance, industrial thermography, autonomous vehicular night vision, and radiometric non-destructive evaluation. To rigorously evaluate the practical feasibility of our proposed two-stage framework, we conduct comprehensive runtime benchmarks across all six representative infrared dataset categories:
\begin{enumerate}
    \item \textbf{Yohann} (Real-world lab calibration scenes, Yohann et al., 2019)
    \item \textbf{DLS-NUC} (Real-world uncooled microbolometer sensor captures, He et al., 2019)
    \item \textbf{Simu w/ TNO} (Simulated degradations over TNO Military Image Fusion scenes)
    \item \textbf{Simu w/ M3FD} (Multi-spectral and multi-sensor aerial targets, Liu et al., 2022)
    \item \textbf{Simu w/ LLVIP} (Low-light and thermal urban pedestrian sequences, Jia et al., 2021)
    \item \textbf{Simu w/ MSRS} (Multi-spectral street road scenes, Tang et al., 2022)
\end{enumerate}

\input{tables/table2_runtime.tex}

\subsubsection{Comparative Runtime Analysis (Table~\ref{tab:runtime_comparison})}
Table~\ref{tab:runtime_comparison} tabulates the execution latency across all 11 baseline algorithms alongside our single-threaded (\textit{Ours}) and multi-core parallelized (\textit{Ours (Parfor)}) implementations on standard $256 \times 256$ frames. 
Filtering baselines such as 1D-GF and 1D-WLS execute with millisecond latency ($\approx 0.003$\,s and $\approx 0.015$\,s) due to their strictly 1D column-wise formulation; however, as demonstrated in Section~\ref{sec:experiments}, they exhibit substantial residual ringing and severe high-frequency edge ghosting. Deep learning baselines (SNRCNN, DLS-NUC, SNRWDNN) execute inference in $0.04-0.12$\,s on modern hardware, yet struggle with out-of-distribution real thermal artifacts without costly sensor-specific retraining. Optimization methods such as LRSID and GSRSR require complex non-convex iterative loops taking between $2.8$\,s and $4.5$\,s.

Our complete two-stage pipeline in single-threaded mode achieves an average per-frame processing latency of approximately $3.55$\,s. When multi-core parallelization is engaged via task-decoupled patch grouping (\textit{Ours (Parfor)}), the execution latency drops to $1.48$\,s across all dataset categories, providing a decisive speedup exceeding $2.3\times$ while simultaneously maintaining state-of-the-art correction quality across all seven quantitative evaluation metrics (Table~\ref{tab:quantitative_evaluation}).

\begin{figure*}[t]
\centering
\includegraphics[width=\textwidth]{figures/fig23_stage_time_comparison.png}
\caption{Stage runtime latency breakdown across the six infrared dataset categories (Section~\ref{sec:runtime_analysis}, Fig.~23). (a) Single-Thread execution mode confirming that Stage~1 (ADMM Destriping) accounts for $\approx 70\%$ of the overall computation budget while Stage~2 (Patch Low-Rank SVD) accounts for $\approx 30\%$. (b) Multi-core parallel execution mode (\textit{Parfor}) demonstrating significant acceleration of the non-local patch grouping and SVD stages, achieving an overall pipeline speedup exceeding $2.3\times$.}
\label{fig:fig23_stage_time_comparison}
\end{figure*}

\subsubsection{Stage Latency Breakdown: 70\% vs. 30\% Ratio (Fig.~23)}
To uncover the internal computational bottleneck of our pipeline, Fig.~\ref{fig:fig23_stage_time_comparison} illustrates the wall-clock execution time decomposed across Stage~1 and Stage~2:
\begin{itemize}
    \item \textbf{Stage~1 (ADMM High-Order Directional Destriping):} Stage~1 iteratively solves an alternating minimization subproblem over $K = 250$ iterations. Per iteration, updating the clean image estimate $\mathbf{C}$ requires 2D Fast Fourier Transforms (FFT/IFFT) costing $\mathcal{O}(M N \log(M N))$, while updating the stripe component $\mathbf{S}$ requires singular value thresholding costing $\mathcal{O}(\min(M^2 N, M N^2))$. On standard $256 \times 256$ frames, Stage~1 consumes $\approx 2.45-2.61$\,s in single-threaded mode, representing $\approx 70.8\%$ of the total runtime budget.
    \item \textbf{Stage~2 (Non-Local Patch Low-Rank SVD):} Stage~2 extracts overlapping $11 \times 11$ patches, matches the top $L = 50$ mutually similar patches across local search windows $\Omega$, constructs observation matrices $\mathbf{G} \in \mathbb{R}^{121 \times 50}$, and computes adaptive singular value thresholding via Eq.~(26). In single-threaded mode, Stage~2 takes $\approx 1.04-1.12$\,s, representing $\approx 29.2\%$ of the total latency.
\end{itemize}

\subsubsection{Multi-Core Parallel Acceleration (Parfor)}
Because reference patch clusters in Stage~2 are statistically decoupled—each group's patch matching, observation matrix assembly, and SVD truncation depend solely on the input image $\mathbf{C}$ and not on adjacent patch groups—Stage~2 represents an embarrassingly parallel computational workload. Distributing these patch clusters across multiple CPU cores via parallel threading (equivalent to MATLAB's \texttt{parfor} or Python's \texttt{joblib}) reduces Stage~2 latency from $\approx 1.08$\,s down to $\approx 0.30$\,s, representing a $> 3.5\times$ speedup for Stage~2 and accelerating the end-to-end framework to $1.48$\,s. 

Applying Amdahl's Law:
\begin{equation}
    S(p) = \frac{1}{(1 - P) + \frac{P}{p}}
\end{equation}
with an empirical parallelizable fraction $P \approx 0.96$ confirms that over $96\%$ of Stage~2's computation executes concurrently, with only $4\%$ attributed to serial memory allocation and final Kaiser window normalization.

% ==============================================================================
% SECTION V: CONCLUSIONS AND FUTURE PROSPECTS
% ==============================================================================
\section{Conclusions and Future Prospects}
\label{sec:conclusions_future_prospects}

\subsection{Summary of Contributions}
In this paper, we presented a novel, robust multistage nonuniformity correction pipeline for single-frame infrared images by harmonizing hybrid high-order directional differential operators and low-rank matrix recovery priors. 
\begin{enumerate}
    \item In Stage~1, an ADMM-based solver effectively separates severe directional stripe non-uniformity $\mathbf{S}$ from the underlying thermal scene $\mathbf{C}$ by coupling a Schatten-$1/2$ quasi-norm low-rank prior on stripes with second-order horizontal directional Total Variation $\left|\frac{\partial^2 \mathbf{C}}{\partial x^2}\right|$.
    \item In Stage~2, residual isotropic environmental noise is eliminated via non-local $11 \times 11$ patch grouping, observation matrix low-rank SVD factorization, and adaptive singular value energy thresholding, followed by Kaiser-windowed back-projection to eliminate blocking artifacts.
    \item Comprehensive quantitative evaluations across six benchmark datasets spanning seven diverse metrics (PSNR, SSIM, EPI, NR-NIQE, BRISQUE, MAE, UIQI) demonstrate that the proposed framework consistently outperforms eleven state-of-the-art filtering, optimization, and deep learning baselines.
\end{enumerate}

\subsection{Future Research Directions & Embedded Hardware Prospects}
Looking forward, several high-impact research trajectories and hardware acceleration opportunities emerge:

\subsubsection{FPGA and DSP Hardware Pipeline for 60+ FPS Video Throughput}
While the current multi-core CPU implementation processes single frames in $\approx 1.48$\,s, modern thermal imaging gimbals, uncrewed aerial vehicles (UAVs), and autonomous drive pods demand video frame rates exceeding $30-60$\,FPS. Because Stage~1 FFTs and matrix additions map directly onto Xilinx/Intel DSP slice pipelines and cuFFT GPU cores, pipelined hardware architectures can stream 2D FFT inversions within clock cycles. Furthermore, Stage~2 patch Euclidean distance matching and $121 \times 50$ SVD operations can be implemented as systolic array cores on FPGA or Tensor Core GPU primitives, paving the way for embedded, sub-$20$\,ms end-to-end video destriping at low electrical power budgets.

\subsubsection{3D Spatio-Temporal Tensor Formulations for Video Infrared Sequences}
For video streams, temporal redundancy between successive frames provides powerful additional prior constraints. By generalizing the current 2D matrix formulation to a 3D spatio-temporal tensor $\boldsymbol{\mathcal{D}} \in \mathbb{R}^{M \times N \times T}$, fixed pattern noise stripes can be modeled as low-tubal-rank tensor structures invariant across the time dimension $T$. Coupling temporal smoothness priors $\left\|\nabla_t \boldsymbol{\mathcal{C}}\right\|_1$ will prevent inter-frame flicker, ghosting, and radiometric drift in continuous video surveillance.

\subsubsection{Self-Supervised and Zero-Shot Physics-Guided Deep Adaptation}
Another promising direction is embedding the exact physical forward model (Eq.~5) and high-order directional priors into deep unfolding networks (such as Learned ISTA or deep ADMM nets). By unrolling $K = 15$ iterations into neural network layers parameterized by learned proximal shrinkages, the model can execute zero-shot physics-guided destriping in single-digit milliseconds without requiring massive paired simulated training corpora.
